\section{Minimax 与可容许性比较整条风险曲线}
\label{sec:11-Mathematical-Statistics/07-Statistical-Decision-and-Reliable-Prediction/03}

两个候选定价规则跑完离线回放，评审桌上摊着两组数。按过去一年的客群构成加权，规则 A 的期望损失比 B 低 \(11\%\)；把客群拆开看，A 在占比 \(3\%\) 的那个新渠道上明显更差，差到那个渠道单独算会亏。有人说看总账，有人说不能让一个渠道兜底。争论持续了半小时，双方谁也没算错。

两边读的是同一条风险曲线的不同部位。上一节按先验把曲线压成一个加权平均，这一节换成盯住曲线的最高点，并给出判断一条规则该不该留在候选名单上的最低门槛。

\subsection{Minimax 盯住曲线的最高点}

若没有可辩护的先验，一种替代是控制最坏状态：

\[
\delta^\ast\in\argmin_\delta\ \sup_{\theta\in\Theta}R(\theta,\delta).
\]

这条规则叫 minimax rule，它承诺的是``无论真实状态落在哪里，风险都不超过某个值''。它不承诺常见状态下表现好，事实上通常反过来——为了压低最高点，容易的区域要让出一些。

拿一个能算到底的例子。观测 \(X\sim\mathcal N(\theta,\sigma^2)\)，平方损失估计 \(\theta\)。朴素规则 \(\delta_0(X)=X\) 的风险恒为 \(\sigma^2\)，曲线是一条水平线。收缩规则 \(\delta_c(X)=cX\)（\(0<c<1\)）的风险为

\[
R(\theta,\delta_c)=(1-c)^2\theta^2+c^2\sigma^2 ,
\]

一条开口向上的抛物线：\(\theta\) 靠近零时方差项的节省占上风，\(\abs{\theta}\) 大了偏差项接管。参数空间若是整条实线，收缩规则的最坏风险无界，minimax 意义下它输给 \(\delta_0\)；若先验知识给出 \(\abs\theta\le B\)，同一条收缩规则在这个有界空间上的最坏风险可能反而更小。同一个规则的稳健性随参数集合改变，所以 minimax 的结论必须连着``上确界是对哪些状态取的''一起读。

\begin{center}
\begin{tikzpicture}[font=\footnotesize,>=stealth,x=1.18cm,y=0.95cm]
  \draw[->,black!70] (-0.15,0) -- (5.5,0) node[below] {\(\theta\)};
  \draw[->,black!70] (0,-0.12) -- (0,4.2) node[above] {风险};
  \draw[black!85,thick] plot[domain=0:5,samples=60] (\x,{0.3+0.13*\x*\x});
  \draw[black!85,thick,dashed] plot[domain=0:5,samples=2] (\x,{1.6});
  \draw[black!40,dotted] (0,3.55) -- (5,3.55);
  \node[anchor=west,black!85] at (5.05,3.55) {\(\delta_A\)};
  \node[anchor=west,black!85] at (5.05,1.6) {\(\delta_B\)};
  \node[anchor=west,black!60] at (0.55,3.85) {\(\sup_\theta R(\theta,\delta_A)=3.55\)};
  \node[anchor=west,black!60] at (0.35,1.85) {\(\sup_\theta R(\theta,\delta_B)=1.60\)};
  \draw[black!30] (0,-1.35) -- (5,-1.35);
  \draw[black!60] plot[domain=0:5,samples=80] (\x,{-1.35+0.9*exp(-2*(\x-1)*(\x-1))});
  \node[anchor=west,black!60] at (5.05,-1.0) {先验 \(\pi\)};
\end{tikzpicture}

\emph{同一对曲线，两种读法：minimax 只看最高点，选 \(\delta_B\)；按下方先验加权平均，\(\delta_A\) 约 \(0.46\) 而 \(\delta_B\) 是 \(1.60\)，Bayes 选 \(\delta_A\)。}
\end{center}

图里两条曲线交叉在 \(\theta\approx 3.1\)。左半边先验质量密集，\(\delta_A\) 省下的是大量常见情形里的代价；右半边先验几乎为零，\(\delta_A\) 在那里的高风险被加权平均忽略不计，却原封不动地留在最高点上。评审桌上那两组数就是这两处：总账读的是曲线与先验的内积，渠道账读的是曲线在某一段上的取值。哪个该赢，取决于新渠道会不会长大，以及那个渠道亏损是否可以被别处的盈利抵掉。

\subsection{可容许性只是一道最低门槛}

若不存在另一条规则在所有 \(\theta\) 上风险都不更大、且至少一点严格更小，就称 \(\delta\) 是可容许的。反过来，被别人整条压住的规则叫不可容许，上一节图里的 \(\delta_3\) 就是。

这道门槛很低。它不保证最大风险小，不保证平均风险小，也不保证现实代价合理；一堆彼此交叉、互不可比的规则可以全部可容许。它只排除一件事：白白多付代价。淘汰完之后要从剩下的里面挑，还得回到先验、最坏情形、对称性之类的准则。

一个反直觉的边界是 James--Stein 现象。设 \(X\sim\mathcal N_d(\theta,I_d)\)，用整个向量的平方欧氏损失评价。逐坐标使用 \(\delta(X)=X\)，每一维单看都无可指摘，但在 \(d\ge 3\) 时这条规则不可容许：适度向原点收缩的 James--Stein 估计量在所有 \(\theta\) 上都把总风险降下来。改进来自允许各维承担一点偏差，换取整体方差的更大下降，而``向原点''这件事依赖于选了哪个原点、用了哪个损失。把多个坐标合进同一个损失以后，逐坐标的最优性不再传递上去。

\subsection{Bayes 与 minimax 由最不利先验相连}

对任意先验 \(\pi\)，平均不超过最大：

\[
\inf_\delta r(\pi,\delta)\ \le\ \inf_\delta\sup_\theta R(\theta,\delta).
\]

于是每个先验都给 minimax 值提供一个下界。若某个先验把这个下界顶到了等号——它的 Bayes rule 的最大风险恰好等于自己的 Bayes 风险——那么这条 Bayes rule 同时是 minimax，这个先验称为最不利先验。有限问题、以及适当紧致且凸的问题里这套论证成立；一般的无限维或非紧问题需要额外正则条件，不能只凭交换 \(\inf\) 与 \(\sup\) 就下结论。

这条联系顺带说明先验还有第二种身份：它可以不代表任何人的信念，只作为构造最坏情形证书的工具。写报告时把这两种身份分开说明，读者才知道该不该对那个先验提意见。

\subsection{平均表现与最坏情形分别对应什么场景}

判断该用哪一边，看决策重复多少次、以及失败能不能摊平。

平台每天做一千万次投放决策，单次损失有界，总代价近似就是期望乘次数，按现实频率加权的平均风险直接对应季度报表，这时 Bayes 视角合适。反过来，一年做两次的产能扩建、错一次就要向监管解释的合规判定、失效不可逆的安全策略，这些决策没有大数定律替你兜底，平均值在其中不代表任何一次实际经历，最坏情形保证才是能承诺出去的东西。

中间地带是子群。把 \(\Theta\) 换成人群划分，最坏情形风险就成了最差子群的风险，这正是公平性评估里那个 worst-group 指标；分布鲁棒优化更进一步，把上确界取在真实分布的某个邻域上，得到的目标介于平均与最坏之间，邻域半径就是保守程度的旋钮。这类目标最后交给的仍是一个带约束的极小极大问题，见\hyperref[sec:09-Convex-Optimization/01-Optimization-Foundations/01]{``把现实选择翻译成优化问题''}里关于目标与约束怎样写全的讨论。

一条实践建议：报告里同时给出加权平均风险和最差子群风险，两个数一起看。只报前者，新渠道那样的问题要等到线上才暴露；只报后者，规则会为了一个可能永远不会长大的小群体而整体保守。

\subsection{Minimax 不是对未知的万能保险}

最坏情形集合选得太大，最优规则会保守到没有信息量——一个永远输出``不确定''的规则在任何状态下风险都不高，也没有任何用处。选得太小，实际会发生的部署变化落在集合外，承诺随之作废。分布鲁棒优化把这件事参数化了，换来的负担是还得选距离、选半径、选损失，每一项都是新的建模决定。

更要紧的是集合本身覆盖不了的东西：标签定义在中途改了口径，训练数据经过了隐藏的筛选，上游机制被自己的决策反过来改变。这些不是分布在邻域里挪了挪，是问题换了一个。稳健性永远要问``对什么扰动稳健''，minimax 给出的是一种清楚的回答格式，不负责替你写出可信的状态空间。

以上三节都假定损失能写下来、概率是可信的。第二个假定该被检验一下：模型报出的 \(0.7\) 到底意味着什么，有没有办法逼它说真话。下一节处理这个问题。
